\chapter{METHODOLOGY}

\section{System Architecture: Machine Learning Pipeline}
The methodology follows a robust machine learning process logic: Data Ingestion $\rightarrow$ Scientific EDA $\rightarrow$ Feature Engineering $\rightarrow$ Hyperparameter Tuning $\rightarrow$ Inference. This modular workflow ensures that each stage is isolated and verifiable (see Appendix~\ref{app:diagrams}, Figure~\ref{fig:ml_logic}).

\section{EDA Rationale and Findings}
Exploratory Data Analysis (EDA) was conducted to scientifically validate the data prior to modeling.
\begin{itemize}
    \item \textbf{Why:} To identify non-stationarity, kurtosis (fat tails), and volatility clustering, which are prevalent in financial data.
    \item \textbf{Expectation:} We expected typical market price series to be non-stationary (I(1) process).
    \item \textbf{Results:} Actual findings from the ADF tests (Table~\ref{tab:adf_results}) confirmed the presence of unit roots. Furthermore, correlation analysis identified that sentiment polarity has a stronger linear relationship with volatility (High-Low spread) than with direct price movement, justifying its use as a regime-filter feature.
\end{itemize}

\section{Scientific Feature Engineering}
Features were engineered to transform raw, non-stationary signals into predictive, scale-invariant inputs:
\begin{itemize}
    \item \textbf{Normalized Technical Distances:} Moving Averages (10, 50) and Bollinger Bands were transformed into percentage distance metrics (e.g., $(Price - MA)/MA$). This ensures the features are stationary and scale-invariant, preventing the model from "learning" absolute price levels which do not generalize across time.
    \item \textbf{Returns-based Momentum:} Differencing the Close price into percentage returns avoids spurious regressions. RSI and MACD (12, 26, 9) quantify the velocity of these returns.
    \item \textbf{Behavioral Signal Fusion:} Daily headlines were converted into a numeric polarity score (-1 to 1) using **FinBERT**. The rationale is based on the transformer's **Self-Attention Mechanism** \citep{vaswani2017}, which weighs contextual tokens to decode sentiment intensity.
\end{itemize}

\section{Market Friction and Backtesting Realism}
A common pitfall in financial ML is the "frictionless assumption." To ensure scientific trust, our backtesting methodology incorporates a **10 basis point (0.1\%) transaction cost** per trade. This accounts for broker commissions and market slippage, ensuring that the reported Sharpe Ratios represent executable performance rather than theoretical artifacts.

\section{Model Implementation: Ensemble Theory}
The prediction task is framed as a supervised binary classification problem. We implement the following models:
\begin{itemize}
    \item \textbf{XGBoost:} A gradient boosting framework that builds decision trees sequentially to minimize a loss function \citep{chen2016}. Theoretically, this reduces the \textbf{Bias-Variance Trade-off} by focusing on the residual errors of previous iterations.
    \item \textbf{Random Forest:} Used as an ensemble benchmark to evaluate performance stability across different asset portfolios.
\end{itemize}

\section{Hyperparameter Optimization}
To prevent underfitting and over-fitting, a grid search was performed over the follow space: \texttt{max\_depth} [3, 5, 7], \texttt{learning\_rate} [0.01, 0.1, 0.2], and \texttt{n\_estimators} [50, 100, 200]. Optimal parameters were selected based on the F1-score to maintain a scientific balance between Precision and Recall.